% Figure: the spectral mapping theorem. A polynomial p(lambda) sends
% each eigenvalue of A to an eigenvalue of p(A). The left panel shows
% three eigenvalues of A on the real axis; the right panel shows their
% images under p(lambda) = lambda^2 - 3 lambda. The eigenvectors are
% carried unchanged, so the arrows connect an input eigenvalue to its
% image, not to a new direction.
\begin{figure}[htbp]
\centering
\begin{tikzpicture}
\begin{axis}[
  width=13cm,
  height=6.6cm,
  xlabel={eigenvalue $\lambda$ of $\mat{A}$},
  ylabel={image $p(\lambda)$},
  xmin=-1.2, xmax=5.2,
  ymin=-3.0, ymax=11.5,
  axis lines=middle,
  tick label style={font=\footnotesize},
  label style={font=\small},
  legend style={font=\footnotesize, at={(0.03,0.97)}, anchor=north west},
  title={Spectral mapping: $p(\lambda)=\lambda^{2}-3\lambda$ acts on the spectrum},
  title style={font=\bfseries},
]
% The polynomial curve.
\addplot[engblue, very thick, domain=-1.1:4.3, samples=80] {x*x - 3*x};
\addlegendentry{$p(\lambda)=\lambda^{2}-3\lambda$}
% Eigenvalues of A on the axis: lambda = 1, 2, 4 (dots at y=0).
\addplot[only marks, mark=*, engred, mark size=2.4pt] coordinates {
  (1,0) (2,0) (4,0)
};
\addlegendentry{eigenvalues of $\mat{A}$}
% Images p(1)=-2, p(2)=-2, p(4)=4 on the image axis (dots at x showing value).
\addplot[only marks, mark=square*, engorange, mark size=2.6pt] coordinates {
  (1,-2) (2,-2) (4,4)
};
\addlegendentry{eigenvalues of $p(\mat{A})$}
% Dashed guides from each eigenvalue up to the curve and across.
\draw[enggray, dashed] (axis cs:1,0) -- (axis cs:1,-2);
\draw[enggray, dashed] (axis cs:2,0) -- (axis cs:2,-2);
\draw[enggray, dashed] (axis cs:4,0) -- (axis cs:4,4);
\node[font=\footnotesize, engorange, anchor=north] at (axis cs:1.5,-2.2)
  {$p(1)=p(2)=-2$};
\node[font=\footnotesize, engorange, anchor=south east] at (axis cs:3.95,4.2)
  {$p(4)=4$};
\end{axis}
\end{tikzpicture}
\caption[The spectral mapping theorem]{The spectral mapping theorem
made visible for $p(\lambda)=\lambda^{2}-3\lambda$ acting on a matrix
$\mat{A}$ with eigenvalues $1$, $2$, and $4$ (red dots on the axis).
The eigenvalues of $p(\mat{A})$ are the images $p(1)=-2$, $p(2)=-2$,
$p(4)=4$ (orange squares), read off the curve. Two distinct eigenvalues
of $\mat{A}$ map to the same eigenvalue $-2$ of $p(\mat{A})$, so a
polynomial can merge a spectrum even when it started well separated;
the two eigenvectors survive the merge and $p(\mat{A})$ acquires a
two-dimensional eigenspace at $-2$. The eigenvectors are carried
unchanged, which is why a function of a matrix is diagonal in the same
basis that diagonalises the matrix.}
\label{fig:vol02:ch10:spectral-mapping}
\end{figure}
